\documentclass[aspectratio=169]{beamer}

\usetheme{Madrid}
\usecolortheme{default}

\usepackage[T1]{fontenc}
\usepackage[utf8]{inputenc}
\usepackage[brazil]{babel}
\usepackage{amsmath}
\usepackage{booktabs}
\usepackage{array}
\usepackage{tikz}
\usetikzlibrary{positioning,arrows.meta,shapes.geometric}
\usepackage{qrcode}

\title[Planejamento da RSL]{Planejamento e Execução de uma Revisão Sistemática}
\subtitle{Pipeline reprodutível: APIs $\to$ Elasticsearch $\to$ LLMs\\DOUTORAMENTO EM INFORMÁTICA (UTAD)}
\author{Alexandre T. Nogueira}
\date{}

\begin{document}

\begin{frame}
\titlepage
\vspace{-1.2em}
\begin{center}
\qrcode[hyperlink,height=2cm]{https://github.com/gameguild-gg/docdo-paper}
\end{center}
\end{frame}

\begin{frame}{Roteiro}
\begin{columns}[t]
\column{0.48\textwidth}
\tableofcontents[hideallsubsections,sections={1-6}]
\column{0.48\textwidth}
\tableofcontents[hideallsubsections,sections={7-12}]
\end{columns}
\end{frame}

% ---------------------------------------------------------------
\section{Visão Geral do Pipeline}
\begin{frame}{1. Visão geral do pipeline}
\centering\vspace{-0.2em}
\begin{tikzpicture}[
	node distance=0.25cm and 0.18cm,
	stage/.style={rectangle, draw, rounded corners, align=center,
	              minimum width=1.75cm, minimum height=0.75cm,
	              inner sep=2pt, fill=blue!8, font=\small},
	tool/.style={rectangle, draw, dashed, align=center,
	             minimum width=1.75cm, minimum height=0.55cm,
	             inner sep=2pt, fill=gray!10, font=\tiny},
	arr/.style={-{Stealth[length=2mm]}, thick},
	cnt/.style={font=\tiny\bfseries, text=blue!70!black}
]
	% Row 1 (left): S1 -> Dedup -> S2a -> S2b
	\node[stage] (s1)  {Coleta};
	\node[stage, right=of s1]  (dd)  {Dedup};
	\node[stage, right=of dd]  (es)  {Elasticsearch};
	\node[stage, right=of es]  (llm) {LLM};

	% Row 1 (right): PDF -> S3 -> QA  (same y as s1 group, placed right of llm)
	\node[stage, right=of llm] (pdf) {PDFs};
	\node[stage, right=of pdf] (s3)  {Full-text};
	\node[stage, right=of s3]  (qa)  {QA};

	% Tool labels (row 2)
	\node[tool, below=of s1]  (t1) {Python\\APIs REST};
	\node[tool, below=of dd]  (t2) {DOI+título};
	\node[tool, below=of es]  (t3) {Título+Resumo};
	\node[tool, below=of llm] (t4) {GPT-4o-mini\\×3 votos};
	\node[tool, below=of pdf] (t5) {Unpaywall\\+manual};
	\node[tool, below=of s3]  (t6) {GPT-5.2\\batch};
	\node[tool, below=of qa]  (t7) {Rubrica\\15Q×3};

	% Counts (row 3)
	\node[cnt, below=0.08cm of t1] {2.985};
	\node[cnt, below=0.08cm of t2] {2.821};
	\node[cnt, below=0.08cm of t3] {638};
	\node[cnt, below=0.08cm of t4] {161};
	\node[cnt, below=0.08cm of t5] {63};
	\node[cnt, below=0.08cm of t6] {52};
	\node[cnt, below=0.08cm of t7] {52};

	% Arrows
	\draw[arr] (s1)--(dd); \draw[arr] (dd)--(es); \draw[arr] (es)--(llm);
	\draw[arr] (llm)--(pdf); \draw[arr] (pdf)--(s3); \draw[arr] (s3)--(qa);
\end{tikzpicture}

\end{frame}

% ---------------------------------------------------------------
\section{Coleta via APIs}
\begin{frame}{2. Coleta de registros via APIs}
\small
Três bases consultadas programaticamente (sem dependência de portais web):
\begin{table}
\centering
\begin{tabular}{>{\raggedright\arraybackslash}p{2.8cm} >{\raggedright\arraybackslash}p{6cm} r}
\toprule
\textbf{Base} & \textbf{API / método} & \textbf{Registros} \\
\midrule
PubMed & Entrez E-utilities (NCBI) --- \texttt{esearch} + \texttt{efetch} XML & 997 \\
arXiv & API oficial do arXiv (Atom XML), categoria \texttt{cs.CV} (Computer Vision) & 904 \\
Semantic Scholar & Graph API v1, paginação por \texttt{offset/limit} & 1.084 \\
\midrule
\textbf{Total} & \textbf{Antes de deduplicação} & \textbf{2.985} \\
\bottomrule
\end{tabular}
\end{table}

\vspace{0.3em}
\textbf{Artefatos gerados:}
\begin{itemize}
	\item \texttt{S1\_search\_results\_REAL.csv} --- 2.985 linhas, metadados brutos.
	\item \texttt{S1\_evidence\_report.md} --- desagregação por fonte, datas de execução.
	\item Script: \texttt{fetch\_all\_real\_data.py} (rate limiting + retry exponencial).
\end{itemize}
\end{frame}

% ---------------------------------------------------------------
\section{Deduplicação}
\begin{frame}{3. Deduplicação}
\small
\textbf{Deduplicação não é trivial:}
\begin{itemize}
	\item Mesmo trabalho aparece em arXiv (preprint) e PubMed (peer-reviewed).
	\item DOIs nem sempre estão presentes em preprints.
	\item Variações no título: caixa, hífens, pontuação, acentos.
\end{itemize}

\textbf{Estratégia em duas chaves:}
\begin{enumerate}
	\item \textbf{DOI exato} (quando disponível): chave primária.
	\item \textbf{Título normalizado}: \texttt{lowercase} $\to$ remover pontuação $\to$ colapsar espaços $\to$ hash.
\end{enumerate}

\begin{table}
\centering
\begin{tabular}{lr}
\toprule
Entrada & 2.985 \\
Duplicatas removidas & 164 \\
\textbf{Saída} (\texttt{S1\_search\_results\_deduplicated.csv}) & \textbf{2.821} \\
\bottomrule
\end{tabular}
\end{table}
\footnotesize Script: \texttt{deduplicate\_s1.py}.
\end{frame}

% ---------------------------------------------------------------
\section{Pré-filtro com Elasticsearch}
\begin{frame}{4a. Pré-filtro ES --- Inclusão (\texttt{bool must})}
\small
\textbf{ES 8.11 (Docker)} --- campo \texttt{combined\_text} = título + resumo; idioma \texttt{english}. \\
\textbf{Todos os 4 grupos obrigatórios} (\texttt{must}), ao menos 1 termo por grupo:

\begin{columns}[T]
\column{0.48\textwidth}
\textbf{DL (IC1):}\\
\scriptsize\texttt{deep learning, neural network,\\convolutional neural, CNN, U-Net,\\UNet, encoder decoder, transformer,\\ResNet, VNet, attention mechanism,\\fully convolutional, FCN, nnU-Net,\\nnUNet, MONAI, autoencoder}

\smallskip
\textbf{Segmentação (IC2):}\\
\scriptsize\texttt{segmentation, segmenting,\\semantic segmentation,\\volumetric segmentation,\\3D segmentation, multi-organ,\\organ segmentation, delineation}

\column{0.48\textwidth}
\textbf{TC/CT (IC3):}\\
\scriptsize\texttt{CT, computed tomography,\\CT scan, CT image, CT volume,\\abdominal CT, chest CT,\\contrast-enhanced, CECT}

\smallskip
\textbf{Órgãos (IC4):}\\
\scriptsize\texttt{liver, kidney, spleen, pancreas,\\lung, heart, stomach, gallbladder,\\bladder, prostate, colon, esophagus,\\adrenal gland, abdominal organ,\\multi-organ, organs at risk, OAR,\\hepatic, renal, pulmonary, cardiac}
\end{columns}

\smallskip\small
\textbf{Pesquisa original com metodologia (IC5)} \textit{(verificado no pós-filtro Python):}\\
\scriptsize
Presenta método, arquitetura ou contribuição metodológica própria.
Exclui: papers só de dataset sem baseline; revisões; notas de aplicação sem metodologia.
\end{frame}

\begin{frame}{4b. Pré-filtro ES --- Exclusões}
\small
\textbf{Nível 1 --- \texttt{must\_not} ES (campo \texttt{title}):}
\quad\texttt{"systematic review"}, \texttt{"literature review"}, \texttt{"survey of"}, \texttt{"a review"}

\vspace{0.3em}
\textbf{Nível 2 --- Pós-filtro Python (\texttt{combined\_text} = título + resumo):}
\begin{table}
\centering\scriptsize
\begin{tabular}{l >{\raggedright\arraybackslash}p{4cm} >{\raggedright\arraybackslash}p{7.5cm}}
\toprule
\textbf{EC} & \textbf{Regra} & \textbf{Condição de exclusão} \\
\midrule
EC1 & Só detecção/classificação & \texttt{detection, classification, diagnosis} presente \textbf{e} \texttt{segment, contour, mask} ausente \\
EC2 & Sem CT no texto & Nenhum de: \texttt{ct\textvisiblespace, computed tomography, ct scan, ct image, ct volume, ct data} \\
EC2 & Modalidade não-TC exclusiva & MRI/US/RX/PET presente \textbf{e} CT ausente \\
EC3 & Alvo não-órgão exclusivo & Tumor/lesão, vaso, osso, músculo, via aérea, COVID --- sem termo de órgão \\
EC4 & Sem 3D/volumétrico & Nenhum de: \texttt{3d, 3-d, three-dimensional, volumetric, volume, v-net, 3d u-net, 3d convolution, 3d cnn} \\
EC5 & Sem deep learning & \texttt{random forest, svm, atlas-based, level set, region growing} sem termos DL \\
EC6 & Revisão/survey & \texttt{systematic review, meta-analysis, overview of, scoping review} no título \\
\bottomrule
\end{tabular}
\end{table}

\centering\small
\textbf{Resultado: 2.821 $\to$ 638} \quad (\texttt{S2\_elasticsearch\_filtered.csv})
\end{frame}

% ---------------------------------------------------------------
\section{Triagem LLM (Título+Resumo)}
\begin{frame}{5. Triagem LLM em cascata}
\small
\textbf{Protocolo de consenso para 638 candidatos:}

\begin{table}
\centering
\begin{tabular}{>{\raggedright\arraybackslash}p{2.4cm} >{\raggedright\arraybackslash}p{2.8cm} c c >{\raggedright\arraybackslash}p{4.5cm}}
\toprule
\textbf{Estágio} & \textbf{Modelo} & \textbf{Runs} & \textbf{Temperature} & \textbf{Função} \\
\midrule
Primário & GPT-4o-mini & 3 & 0.3 & Triagem inicial; \textbf{INCLUDE unânime} obrigatório \\
Validação & GPT-5-nano & 3 & default & Comparação independente \\
Desempate & GPT-5.2 & 1 & 0.3 & Resolver disagreement (274 papers) \\
\bottomrule
\end{tabular}
\end{table}

\textbf{Prompt estruturado:} JSON output com \texttt{decision}, IC1--IC6, EC1--EC6, \texttt{confidence}, \texttt{reasoning}.

\textbf{Regra conservadora:} qualquer divergência entre as 3 chamadas $\to$ desempatador GPT-5.2.

\centering
\textbf{638 $\to$ 161 incluídos} \quad (\texttt{all\_papers\_after\_s2\_with\_status.csv})
\end{frame}

% ---------------------------------------------------------------
\section{Validação Humana}
\begin{frame}{6. Validação humana da triagem LLM}
\small
\textbf{Auditoria 10\%}, amostragem estratificada (n = 64 dos 638 pós-Elasticsearch).

\textbf{Estratos:} decisão da IA (INCLUDE/EXCLUDE) $\times$ tipo de consenso (unânime/maioria) $\times$ base de origem.

\begin{table}
\centering
\begin{tabular}{lr}
\toprule
Concordância IA--humano & \textbf{96,4\%} (62/64) \\
Cohen's $\kappa$ & \textbf{0,89} \\
Falsos negativos (IA EXCLUI $\to$ humano INCLUI) & 3 \\
Falsos positivos (IA INCLUI $\to$ humano EXCLUI) & 0 \\
\bottomrule
\end{tabular}
\end{table}
\footnotesize Os 3 falsos negativos foram reintroduzidos na triagem de texto completo. Relatório: \texttt{S6\_validation\_report.md}.
\end{frame}

% ---------------------------------------------------------------
\section{Aquisição de PDFs}
\begin{frame}{7. Aquisição de texto completo}
\small
Dos 161 estudos selecionados em S2b, foi necessário obter o PDF para a triagem de texto completo.

\textbf{Estratégia:}
\begin{enumerate}
	\item Tentativa automática via \textbf{Unpaywall} (DOI $\to$ URL open access).
	\item Captura direta via arXiv (PDF público).
	% \item Acesso institucional (UTAD/sci-hub apenas para verificação de elegibilidade).
	\item Registro de falhas em \texttt{excluded\_no\_fulltext.csv}.
\end{enumerate}

\begin{table}
\centering
\begin{tabular}{lr}
\toprule
Selecionados em S2b & 161 \\
PDFs obtidos & \textbf{63} \\
Sem acesso (paywall ou indisponível) & 98 \\
\bottomrule
\end{tabular}
\end{table}
\footnotesize Detalhamento em \texttt{S11\_paywalled\_access.md}.
\end{frame}

% ---------------------------------------------------------------
\section{Triagem de Texto Completo}
\begin{frame}{8. Triagem de texto completo (LLM batch)}
\small
\textbf{Pipeline:}
\begin{enumerate}
	\item Extração de texto: \texttt{PyMuPDF} $\to$ \texttt{markdown} normalizado.
	\item Truncamento controlado para caber na janela do modelo.
	\item Submissão em lote (\texttt{batch API}) ao \textbf{GPT-5.2}.
	\item Re-aplicação dos critérios IC1--IC6 / EC1--EC6 sobre o texto completo.
	\item Extração estruturada: dataset, métricas (Dice/HD), arquitetura, órgãos, dataset, GPU.
\end{enumerate}

\begin{table}
\centering
\begin{tabular}{lr}
\toprule
PDFs avaliados & 63 \\
\textbf{Incluídos na síntese} & \textbf{52} \\
Excluídos (motivos categorizados) & 11 \\
\bottomrule
\end{tabular}
\end{table}

\footnotesize Saídas: \texttt{s3\_extracted\_data\_full.csv} (52 linhas, dados extraídos), \texttt{s3\_excluded\_papers.csv} (motivos), \texttt{s3\_summary\_table.csv}.
\end{frame}

% ---------------------------------------------------------------
\section{Quality Assessment}
\begin{frame}{9. Quality Assessment (rubrica)}
\small
\textbf{Rubrica de 15 questões $\times$ 3 dimensões} (5 questões cada, escala 0--2 pontos; total 0--30):

\begin{itemize}
	\item \textbf{Rigor metodológico:} clareza do método, baselines, hiperparâmetros, código disponível, ablações.
	\item \textbf{Reprodutibilidade:} dataset público, código público, seeds, métricas padronizadas, especificação de hardware.
	\item \textbf{Relevância clínica:} validação externa, comparação com radiologistas, métricas clínicas, tempo de inferência, integração de fluxo.
\end{itemize}

\textbf{Aplicação:} GPT-5.2 batch sobre os 52 incluídos.

\begin{table}
\centering
\begin{tabular}{lr}
\toprule
Pontuação média total & \textbf{21,0 / 30} \\
Amplitude & 14 -- 28 \\
Alta qualidade ($\geq 24$) & 10 (19,2\%) \\
Média qualidade (15--23) & 41 (78,8\%) \\
Baixa qualidade ($<15$) & 1 (1,9\%) \\
\bottomrule
\end{tabular}
\end{table}
\footnotesize Per-study scores: \texttt{qa\_summary.csv}; respostas brutas: \texttt{qa\_parsed\_results.json}.
\end{frame}

% ---------------------------------------------------------------
\section{Artefatos e Rastreabilidade}
\begin{frame}{10. Artefatos suplementares (S0--S12)}
\small
Cada etapa do pipeline gera CSVs/MDs versionados que permitem auditoria:

\begin{columns}[T]
\column{0.5\textwidth}
\textbf{Protocolo e critérios:}
\begin{itemize}
	\item \texttt{S0\_data\_provenance.md}
	\item \texttt{S3\_search\_protocol.md}
	\item \texttt{S4\_ai\_screening\_protocol.md}
	\item \texttt{S4\_screening\_criteria.md}
\end{itemize}

\textbf{Dados brutos e filtrados:}
\begin{itemize}
	\item \texttt{S1\_search\_results\_REAL.csv} (2.985)
	\item \texttt{S1\_*\_deduplicated.csv} (2.821)
	\item \texttt{S2\_elasticsearch\_filtered.csv} (638)
\end{itemize}

\column{0.5\textwidth}
\textbf{Decisões e auditoria:}
\begin{itemize}
	\item \texttt{S5\_screening\_decisions.csv}
	\item \texttt{S6\_validation\_report.md}
	\item \texttt{s3\_excluded\_papers.csv}
\end{itemize}

\textbf{Síntese final:}
\begin{itemize}
	\item \texttt{S2\_final\_included\_studies.csv} (52)
	\item \texttt{s3\_extracted\_data\_full.csv}
	\item \texttt{qa\_summary.csv}
	\item \texttt{S12\_per\_organ\_statistics.md}
\end{itemize}
\end{columns}
\end{frame}

% ---------------------------------------------------------------
\section{Funil PRISMA}
\begin{frame}{11. Funil de seleção (PRISMA modificado)}
\centering
\scriptsize
\begin{tikzpicture}[
		node distance=0.45cm,
		box/.style={rectangle, draw, rounded corners, align=center, minimum width=8.3cm, minimum height=0.55cm, inner sep=2.5pt},
		side/.style={rectangle, draw, align=center, minimum width=3.6cm, minimum height=0.55cm, inner sep=2pt},
		arr/.style={-{Stealth[length=2mm]}, thick}
]
	\node[box] (s1raw) {S1: Registros identificados (PubMed + arXiv + Semantic Scholar): \textbf{2.985}};
	\node[box, below=of s1raw] (dedup) {Após deduplicação (DOI + título normalizado): \textbf{2.821}};
	\node[side, right=0.2cm of dedup] (dedupexc) {Duplicatas removidas: \textbf{164}};
	\node[box, below=of dedup] (s2) {S2a: Pré-filtro Elasticsearch (título + resumo): \textbf{638}};
	\node[side, right=0.2cm of s2] (s2exc) {Excluídos no pré-filtro: \textbf{2.183}};
	\node[box, below=of s2] (llm) {S2b: Triagem LLM (GPT-4o-mini ×3 unânime): \textbf{161}};
	\node[side, right=0.2cm of llm] (llmexc) {Excluídos via LLM: \textbf{477}};
	\node[box, below=of llm] (pdf) {Texto completo obtido (PDF): \textbf{63}};
	\node[side, right=0.2cm of pdf] (nopdf) {Sem PDF: \textbf{98}};
	\node[box, below=of pdf] (s3) {S3: Triagem de texto completo (GPT-5.2): \textbf{52 incluídos}};
	\node[side, right=0.2cm of s3] (s3exc) {Excluídos (texto completo): \textbf{11}};
	\node[box, below=of s3] (final) {Síntese narrativa + QA: \textbf{52 estudos}};

	\draw[arr] (s1raw) -- (dedup);
	\draw[arr] (dedup) -- (s2);
	\draw[arr] (s2) -- (llm);
	\draw[arr] (llm) -- (pdf);
	\draw[arr] (pdf) -- (s3);
	\draw[arr] (s3) -- (final);
	\draw[arr] (dedup.east) -- (dedupexc.west);
	\draw[arr] (s2.east) -- (s2exc.west);
	\draw[arr] (llm.east) -- (llmexc.west);
	\draw[arr] (pdf.east) -- (nopdf.west);
	\draw[arr] (s3.east) -- (s3exc.west);
\end{tikzpicture}
\end{frame}

% ---------------------------------------------------------------
\section{Reprodutibilidade}
\begin{frame}{12. Reprodutibilidade e limitações}
\small
\textbf{O que é reprodutível bit-a-bit:}
\begin{itemize}
	\item Etapa de deduplicação (determinística sobre \texttt{S1\_search\_results\_REAL.csv}).
	\item Pré-filtro Elasticsearch (consulta + versão do índice fixadas).
	\item Cálculo das estatísticas e tabelas a partir dos CSVs finais.
\end{itemize}

\textbf{O que \emph{não} é reprodutível bit-a-bit:}
\begin{itemize}
	\item Coleta S1 --- APIs públicas retornam conjuntos diferentes ao longo do tempo.
	\item Decisões LLM --- modelos não são estritamente determinísticos (mesmo com $T=0$).
\end{itemize}

\textbf{Mitigação:} todos os artefatos intermediários estão versionados em \texttt{pipeline\_outputs/}, permitindo auditoria do estado em que cada decisão foi tomada.
\end{frame}

% ---------------------------------------------------------------
\begin{frame}{Próximos passos e materiais}
\small
\begin{columns}[c]
\column{0.65\textwidth}
\begin{itemize}
  \item \textbf{Generalizar} o pipeline como ferramenta reutilizável de RSL assistida por IA.
  \item \textbf{Comparar} com revisores humanos em outro tópico (validação externa).
  \item \textbf{Publicar} o protocolo no PROSPERO e o código sob licença aberta.
  \item \textbf{Integrar} ao fluxo do projeto DocDo (planejamento cirúrgico).
\end{itemize}
\column{0.30\textwidth}
\centering
\qrcode[hyperlink,height=4.5cm]{https://github.com/gameguild-gg/docdo-paper}\\
\scriptsize Repositório e artefatos
\end{columns}
\end{frame}

\end{document}
